\textbf{実行可能な位相．}見出しの許容容量 $\Upk\le2$ で解が得られる離手位相は，往路では $\phi_\ell=0.062$〜$0.25$ が全体重に共通である．
復路で $\Upk\le2$ の解が得られたのは $m=$60，63，72\,kg だけで，位相は $\phi_\ell=0.562$〜$0.938$ である．
$\phi_\ell=0.438$〜$0.5$ では，どの体重でも $\Upk\le2$ の解は得られなかった．B が最も遠い $\phi=0.5$ の前後で，先端間距離は 2.6〜2.7\,m になり，A は最も低い位置にある．片側からの振りでは，この距離を越える離手の速度と高さが得られない．

\textbf{体重と装置周期．}必要把持容量の最小値は，$m=60$\,kg の \capSixty から $m=75$\,kg の \capSeventyFive へ体重とともに増える（体重比では \capBWmin〜\capBWmax\,BW）．関節トルクの上限を体重によらず固定しているので，重い体ほど振りに使える余力が小さく，体重比でも要求が増える．装置周期への依存は小さい．$\phi_\ell=0.25$ では，9 通りの周期の間の $\Upk$ の違いは最大 \periodSpreadMax\,\%（中央値 \periodSpreadMed\,\%）である．装置の速さは 0.08〜0.11\,m/s で，離手のときの体の速さ（数 m/s）に比べて小さい．飛行時間は \flightMin〜\flightMax\,s で，その間に装置が動く距離は 5\,cm 以下である．効くのは離手の瞬間の装置の位置であって，速さではない．

\textbf{最大負荷が生じる相．}$\Upk\le2$ で解けた \gridFeasTwo 条件のうち，保持の最大負荷が $\Upk$ に達しているものは \holdAtPeak\,\%，振りの最大負荷が $\Upk$ に達しているものは \swingAtPeak\,\%である．多くの解で両方が同時に上限に触れている．どちらが必要容量を決めているかは，到達しているかどうかでは分からない．これを第~\ref{sec:interventions}節で調べる．
